\documentclass[12pt]{article}

\usepackage{sbc-template}
\usepackage{graphicx,url}
\usepackage[english]{babel}
\usepackage[utf8]{inputenc}


\sloppy

\title{Proposal for Assignment 2: Statistical Analysis and Demand Prediction for Bike Sharing: A DOE-Based Approach}

\author{[Student 1 Name]\inst{1}, [Student 2 Name]\inst{1}}

\address{[Department/University Name]\\
  [Course Name] -- Design of Experiments
}

\begin{document}

\maketitle

\begin{abstract}
This proposal presents our plan for the second assignment of the Design of Experiments course. We intend to perform a comprehensive statistical analysis on the UCI Bike Sharing dataset, using Design of Experiments (DOE) reasoning to identify factors, levels, effects, and interactions, and subsequently develop a well-justified multiple linear regression model to predict total bike rental demand ($cnt$).
\end{abstract}

\begin{resumo}
Esta proposta apresenta o plano para o segundo trabalho da disciplina de \textit{Design of Experiments}. Pretendemos realizar uma análise estatística abrangente sobre o dataset \textit{Bike Sharing} da UCI, utilizando raciocínio de \textit{Design of Experiments} (DOE) para identificar fatores, níveis, efeitos e interações e, a partir disso, desenvolver um modelo de regressão linear múltipla bem justificado para prever a demanda total de bicicletas ($cnt$).
\end{resumo}

\section{Dataset}

We chose the \textbf{Bike Sharing Dataset} from the UCI Machine Learning Repository (Fanaee-T \& Gama, 2013), collected from the Capital Bikeshare system in Washington D.C., USA, covering the years 2011 and 2012.

The dataset is provided in two aggregated forms: \texttt{day.csv} (731 days) and \texttt{hour.csv} (17,379 hours). For this work, we will use \textbf{\texttt{hour.csv}} since it includes the \texttt{hr} (hour of day) variable, which is fundamental to capture the hourly demand profile and to investigate relevant interactions.

The response variable is \textbf{\texttt{cnt}} = total number of bike rentals (\texttt{casual + registered}). The explanatory variables include temporal factors (\texttt{season}, \texttt{yr}, \texttt{mnth}, \texttt{hr}, \texttt{holiday}, \texttt{weekday}, \texttt{workingday}) and meteorological covariates (\texttt{weathersit}, \texttt{temp}, \texttt{atemp}, \texttt{hum}, \texttt{windspeed}). Variables \texttt{casual} and \texttt{registered} will not be used as predictors, as they are components of the response.

\section{General Objective}

The general objective is to \textbf{perform a DOE-based statistical analysis of bike rental demand} to: (i) identify which factors and interactions significantly affect the total number of rentals (\texttt{cnt}); (ii) understand the behavior of demand under different temporal and weather conditions; and (iii) develop and validate a \textbf{well-justified multiple linear regression model} as a predictor for \texttt{cnt}, supported by effect plots, ANOVA, and regression diagnostics.

\section{Proposed Methods}

Since the data come from a real operating system with no controlled randomization, we treat this as an \textbf{observational study with a natural factorial structure}. Our methodology is structured as follows:

\subsection{Experimental Design Characterization}
\begin{itemize}
\item \textbf{Design type:} Natural factorial design (multiple factors acting simultaneously).
\item \textbf{Unit:} One hour of a calendar day (\texttt{hour.csv}).
\item \textbf{Randomization:} None (real system data).
\item \textbf{Structure:} Primarily between-observations with a temporal/sequential structure, which may introduce autocorrelation.
\end{itemize}

\subsection{Exploratory and Descriptive Analysis}
We will conduct EDA including: sample characterization, summary statistics (mean, median, standard deviation, min, max, quartiles), histograms/density plots, boxplots by \texttt{hr}, \texttt{workingday}, \texttt{weathersit}, and \texttt{season}, time series visualization (2011 vs.\ 2012), and correlation analysis among numerical variables.

\subsection{DOE-Based Statistical Analysis}
We will focus on factors with strong practical/statistical relevance: \texttt{workingday}, \texttt{hr}, \texttt{weathersit}, \texttt{season}, and \texttt{temp}. Our analysis will include:
\begin{itemize}
\item \textbf{Main effects plots} to visualize the magnitude/direction of effects.
\item \textbf{Interaction plots}, with priority to \textbf{\texttt{workingday * hr}} (fundamental due to the distinct commute vs.\ leisure profiles), followed by \texttt{weathersit * hr} and \texttt{season * hr}.
\item \textbf{One-way and factorial ANOVA}, followed by \textbf{Tukey HSD post-hoc tests} when appropriate, to statistically support observed effects and interactions.
\end{itemize}

\subsection{Regression Modeling, Transformations and Diagnostics}
We will build models progressively:
\begin{itemize}
\item \textbf{Modeling strategy:} Start from main effects (M0), add \texttt{workingday * hr} (M1), test additional interactions (M2), assess non-linearity via \texttt{temp$^2$} (M3), and obtain a parsimonious final model (M4).
\item \textbf{Categorical treatment:} Use dummy coding for categorical factors. Highly correlated pairs (e.g., \texttt{temp} vs.\ \texttt{atemp}) will be handled by selecting \texttt{temp} only.
\item \textbf{Transformations:} Evaluate Box-Cox and alternatives (\texttt{sqrt(cnt)}, \texttt{log(cnt+1)}) to improve normality/homoscedasticity if needed, justifying any transformation.
\item \textbf{Assumptions:} Check normality (QQ-plots/Shapiro-Wilk), homoscedasticity (Residuals vs.\ Fitted, Levene/Bartlett), independence/temporal dependence (\textbf{Durbin-Watson}, \textbf{ACF/PACF}), multicollinearity (\textbf{VIF}), and influential points (Cook's Distance/leverage).
\item \textbf{Validation:} Use \textbf{temporal split} (train = 2011, test = 2012) to evaluate generalization. Metrics: $R^2$, adjusted $R^2$, AIC, BIC, RMSE, and MAE. Model selection will balance fit, parsimony, interpretability, and DOE coherence.
\end{itemize}

\section{References}

\small
\noindent [1] Fanaee-T, H. and Gama, J. (2013). Event labeling combining ensemble detectors and background knowledge. \textit{Progress in Artificial Intelligence}, pp. 1--15, Springer Berlin Heidelberg. doi:10.1007/s13748-013-0040-3.

\noindent [2] UCI Machine Learning Repository. Bike Sharing Dataset. Available at: \url{https://archive.ics.uci.edu/ml/datasets/Bike+Sharing+Dataset}.

\end{document}
